\section{A direct quartic in 314 unrestricted integers}
\label{sec:direct314}

This component uses the numerical C222 proof calculus. It has no K215
simulation and no specialization of a universal polynomial. Its complete
ordinary polynomial is
\begin{equation}\label{eq:direct314}
 P_{314}(z_0,\ldots,z_{313})=\sum_{j=0}^{203}r_j(z)^2=0,
 \qquad z\in\Z^{314}.
\end{equation}
Every quadratic residual is printed in the appendix and supplied literally
in \code{direct/equation.json}; every collected coefficient is in
\code{direct/expanded_polynomial.json}. The new equation has exact degree
four, 5,721 monomials and coefficient height 128,446.

\begin{theorem}\label{thm:direct314}
With the retained ordinary C222-to-ZFC refutation calibration,
\[
 \ZFC\vdash
 \bigl(\neg\exists z\in\Z^{314}\;P_{314}(z)=0\bigr)
 \ \longleftrightarrow\ \Con(\ZFC).
\]
\end{theorem}

Here ``retained calibration'' identifies previously proved mathematical
dependencies; it is not an added consistency hypothesis. The present
package does not claim a new end-to-end proof-assistant derivation.
The exact parent comparison, independently recomputed from the literal
equations, is
\begin{center}
\begin{tabular}{lrrrrr}
\toprule
Construction & Variables & Rows & Degree & Terms & Height\\
\midrule
Package 131 parent &324&208&4&5,808&128,446\\
This construction &314&204&4&5,721&128,446\\
\bottomrule
\end{tabular}
\end{center}
Thus the reduction saves ten coordinates, four quadratic rows and 87
collected monomials, without increasing degree or coefficient height.
It improves this direct-calculus comparison; other coefficient
architectures already use fewer variables.

\subsection{The arithmetic interface and the semantic anchor}

The pinned parent is package 131 at commit
\code{1c5367452a0bd4dd7e9e4bceb1508b62dbd2ff61}. Its accepted factored
quadratic source is retained as \code{direct/baseline/factored.json}, and
its complete arithmetic argument as \code{direct/proofs/parent131.md}.
The front has 18 power requests sharing 15 independent Pell pairs.
One requested exponent is the derived syntax-label polynomial
\[
 h=J(A_a,A_b,A_c,A_d).
\]
The corresponding pair is literally \code{index15.X,index15.Y}, and
its exponent is \code{1+add.380}. No new unknown is introduced by
writing $h$.

For the numerical implication constructor
\[
 I(a,b)=2(a+b)(a+b+1)+4b,
\]
$J$ is the fixed sum of the nine axiom conclusions, $a,b,I(a,b)$,
the universal quantification constructor and 1. Its original natural circuit has
nonnegative coefficients. The retained encoded-column guards give
$A_a,A_b,A_c,A_d\ge1$, so
\begin{equation}\label{eq:Janchor}
 h\ge16,\qquad h\ge I(A_a,A_b)\ge2A_a^2.
\end{equation}
These are raw circuit consequences, prior to interpreting any proof.

The parent has an additional fully indexed Pell cap at its frame index
$W$. We reuse the pair $(U,C)$ already requested at $h$ as that cap.
Its principal norm is retained; the old separate cap norm is deleted.
The index-error pool continues to include all fifteen principal pairs,
including the anchor. The new frame and cap conditions are
\begin{align}
 W&=2\left(2+\sum_{a\in\mathcal A}a^2\right),\label{eq:evenframe314}\\
 C^2-W^2-1-\sum_{i\ne *}y_i^2&\ge0.\label{eq:squarecap314}
\end{align}
Here $\mathcal A$ is precisely the retained list of all actual base,
exponent and output atoms, including $h$, with their complete affine
restorations. The asterisk marks the anchor. The second condition is
one quadratic row using
\[
 \mathcal N(a,b,c)=a^2+b^2+c(c+1).
\]
This form is nonnegative on $\Z^3$ and represents every natural integer:
apply the three-square theorem to $4n+1$, whose three square residues
force one odd and two even coordinates.

All other front signs and arithmetic requests are retained. We remove
the old sign guard on the scale $A$ and the old sign guard on the
auxiliary multiplier, with their three private witnesses each. The
ordinary sign argument is given below, rather than borrowing those
deleted hypotheses.

\subsection{An odd, large scale with either sign}

Write $X_n(b),Y_n(b)$ for the integer Pell polynomials, defined by
$X_0=1,Y_0=0$ and
\[
 X_{n+1}=bX_n+(b^2-1)Y_n,
 \qquad Y_{n+1}=X_n+bY_n.
\]
For negative bases they satisfy
\[
 X_n(-b)=(-1)^nX_n(b),\qquad
 Y_n(-b)=(-1)^{n-1}Y_n(b)\quad(n\ge1).
\]
Negative indices use $X_{-n}=X_n$ and $Y_{-n}=-Y_n$.

The original odd-ordinate frame norm remains
\begin{equation}\label{eq:framenorm314}
 A^2-W^3(W+2)(2z+1)^2=1.
\end{equation}
Equation~\eqref{eq:evenframe314} makes $W$ even and $W\ge516$, hence
$A$ is odd. Put $a=|A|$. Classification at center $W+1$ shows that
$(a,|W(2z+1)|)$ has a positive Pell index $n$. As
$Y_n(W+1)\equiv n\pmod W$, this index is a positive multiple of $W$.
Consequently
\begin{equation}\label{eq:scale314}
 a\ge X_W(W+1)>W^W,
 \qquad |A-1|\ge a-1>4W^3.
\end{equation}
Every semantic atom has absolute value below $W$. In particular
$|k_i|<W$ for every requested index and $16\le h<W$.

\subsection{A deterministic comparison center}

The crucial new local construction replaces a CRT choice by the
polynomial
\begin{equation}\label{eq:detcenter314}
 \boxed{D=A+(1-A)H^2.}
\end{equation}
The norms and auxiliary multiplier are
\begin{align}
 U^2-(A^2-1)C^2&=1,\label{eq:anchornorm314}\\
 H^2-(A^2-1)J_c^2&=1,\label{eq:commonnorm314}\\
 J_c&=C^2(2e+1),\qquad e\in\Z.\label{eq:oddaux314}
\end{align}
The cap gives $|C|>W$, so $C,J_c$ are nonzero. From the common norm,
\begin{equation}\label{eq:centeridentity314}
 D=1-(A-1)(A^2-1)J_c^2.
\end{equation}
Since $A$ is odd and $C^2\mid J_c$, equations
\eqref{eq:detcenter314}--\eqref{eq:centeridentity314} imply
\begin{equation}\label{eq:centerresidues314}
 D\equiv A\pmod H,\qquad D\equiv1\pmod{4C},\qquad |D|>1.
\end{equation}
For $A\ge3$, the second expression is less than $-1$; for $A\le-3$
it exceeds 1. Thus neither a sign condition on $D$ nor an adjustable
comparison-center parameter is needed.

Retain the comparison relations
\begin{equation}\label{eq:comparison314}
 S^2-(D^2-1)Q^2=1,
 \qquad S=U+Hg,
 \qquad Q=h+4Cq.
\end{equation}
We now establish the exact absolute anchor index.

\begin{lemma}\label{lem:anchor314}
Under the displayed conditions, $|C|=Y_h(|A|)$.
\end{lemma}
\begin{proof}
Classify the first two norms at $a=|A|$:
\[
 |U|=X_i(a),\quad |C|=Y_i(a),\quad
 |H|=X_l(a),\quad |J_c|=Y_l(a),
\]
with positive indices $i,l$. The divisibility $C^2\mid J_c$ gives
$i<l$ and $|C|\mid l$. The latter is the usual squared-ordinate lemma:
first $Y_i\mid Y_l$ implies $i\mid l$ by the addition law. Write
$l=ik$ and expand the Pell unit power modulo $Y_i^2$ to obtain
\[
 Y_l/Y_i\equiv kX_i^{k-1}\pmod{Y_i}.
\]
The norm gives $\gcd(X_i,Y_i)=1$, so $Y_i\mid k$ and hence $Y_i\mid l$.

Classify the comparison pair at $|D|$. Polynomial dependence on the
base, $D\equiv A\pmod H$, and $S\equiv U\pmod H$, with the independent
coordinate signs absorbed, imply that its index $j$ satisfies
$j\equiv\pm i\pmod{2l}$. This is the parent's abscissa-residue lemma:
reduce the unit exponent modulo $2l$, reflect about $l$, and use strict
growth to distinguish the remaining absolute traces. Signed bases
contribute only further allowed signs.

Because $D\equiv1\pmod{4C}$, one has $Y_j(D)\equiv j\pmod{4C}$.
The last comparison congruence gives $j\equiv\pm h\pmod{4C}$.
As $|C|\mid l$, these imply $i\equiv\pm h\pmod{2|C|}$. But
$1\le i\le|C|$ and $1\le h<W<|C|$, so $i=h$.
\end{proof}

The two quadratic rows
\[
 T-H^2=0,\qquad D-A-(1-A)T=0
\]
realize the deterministic center. They replace the old presentation
$D=1+4J_c(2d+1)$, $D=A+Hf$. Actual savings are measured after full
projection; a local variable subtraction alone would not establish them.

\subsection{Signed index recovery without a sign guard on the scale}

Each principal norm gives
$y_i=\epsilon_iY_{j_i}(A)$ and $x_i=\pm X_{j_i}(A)$ for some
$j_i\ge0$. For non-anchor pairs the square cap implies
$|y_i|<|C|$, hence $j_i<h$; for the anchor $j_i=h$. All are below $W$.
The unchanged pooled row is
\begin{equation}\label{eq:pool314}
 \sum_{i=1}^{15}(y_i-k_i)^2=(A-1)q_0,
 \qquad q_0\in\Z.
\end{equation}
The polynomial recurrence gives $Y_j(A)\equiv j\pmod{A-1}$ for
every signed $A$. Thus $|A-1|$ divides
\[
 E=\sum_{i=1}^{15}(\epsilon_i j_i-k_i)^2.
\]
This finite index error satisfies
$0\le E<4\cdot15W^2<4W^3<|A-1|$. Hence $E=0$, and
\begin{equation}\label{eq:canonical314}
 y_i=Y_{k_i}(A),\qquad x_i=\pm X_{k_i}(A)
\end{equation}
for all requested indices, including zero or negative ones before the
front signs have been decoded. In particular $C=Y_h(A)$ with the correct
signed-base convention. We bound $E$, not the enormous ordinate-error
sum in~\eqref{eq:pool314}.

This signed-base argument credits the repository's earlier unoriented-cap
work, retained as \code{direct/proofs/signed_scale.md}. The present
proof rebinds it to the new complete semantic source.

\subsection{Power semantics and the inherited arithmetic endpoint}

For $k>0$, $2\le n<W$ and $|m|<W$, put $T_n=2An-n^2-1$ and
\[
 L_k=X_k(A)+(n-A)Y_k(A),\qquad
 R_k=X_k(A)-(n-A)Y_k(A).
\]
The identities
\[
 L_kR_k=1+T_nY_k(A)^2,\qquad
 L_{k+1}=nL_k+T_nY_k(A)
\]
hold for either sign of $A$, so $L_k\equiv n^k\pmod{T_n}$. The
frame gives $|T_n|>W^W$, larger than both $|m-n^k|$ and $|mn^k+1|$.
The genuine request therefore gives $m=n^k$ on the canonical abscissa
branch. The opposite branch would give $mn^k=-1$, impossible in integers.
Canonical orientation here need not mean a positive coordinate.

The initial retained rows give
$B=r+s+1\ge1$, $q=B+o+1\ge2$, $x\ge0$ and
$t=16+h_0^2\ge16$. The $q^t$ request first selects the canonical
$t$-pair. Its second request gives $V_0=(F^2+1)^t$, including the
possible base-one case. The ceiling $FB=V_0+s>0$ yields $F\ge1$,
then $B>F$ and $q\ge3$. The cofinal equation
$y=V_0(V_0-1)$ yields $y\ge F+1$. Retained encoded-column guards
then make every requested index positive, including
$K=2(2+t+x+y+F)$. The $B^K$ and $(BH_0)^K$ requests recover
\[
 N=B^K,
 \qquad q=(BH_0)^K,
 \qquad H_0=1+\sum \text{the nine positive representatives}.
\]
All remaining pairs support requests at $q$ or $q+1$, and the polynomial
addition law supplies the derived $t+1$ pair. Thus all 18 power
interfaces are genuine. Each retained interval-product guard, together
with $r_i+s_i=L_i-1\ge0$, then restores both remainder signs.

At this point the parent's ordinary arithmetic theorem applies to the
unchanged front. Its exact proof is retained, with the following chain
made explicit. The ceiling gives $V_0\le BF<B^2\le q$, so the factorial
digit is exact: $B=\binom{F^2+t}{t}$. The parent's fixed-point theorem
at $t\ge16$ yields $F=t!$. The same $B$ equals the product of the
pairwise coprime integers $d_i=1+F^2/i$, $1\le i\le t$. Amplified
prime-power ranges at $N=B^K$, with the self-sized radix $NH_0^K$,
recover every bounded local rule and derived-label constraint. The
Horn-flow multiplicity argument then gives a C222 refutation, even for
a supplied multiset containing cycles. The named logical interfaces
\code{KEndpoint.c222_iff_zf_refutation} and
\code{ConstructibleChoice.zf_refutation_iff_zfc_refutation}
give a finite ZFC refutation. The older full formal development is
identified rather than recompiled in this release.

\subsection{Complete witness construction}

Conversely, begin with a finite C222 refutation. Use the parent's
cofinal horizon $t=16+h_0^2$, truth padding, factorial, beta code and
front construction. Fix $R_0,K,B,N$ and choose the nine positive CRT
representatives modulo $d_i^{2K}$. Their common period is $N^2$.
Increase only $A_a$ by a sufficiently large multiple of $N^2$ so that
\begin{equation}\label{eq:dominate314}
 h\ge2\max_{i\ne *}k_i+1.
\end{equation}
This is possible by~\eqref{eq:Janchor}: the other varying target is
only $A_a$ itself, whereas $h\ge2A_a^2$. Every local polynomial
congruence is preserved at the original stronger precision. The parent
has no aggregate upper bound on the representatives.

Only after this lift compute $H_0$, $q=(BH_0)^K$, all power outputs,
digits and finally the doubled frame $W$. The uniform estimate
$J\le H_0^{2^{22}}$ and $2^{22}(t+1)<K$ preserve all digit bounds;
they do not require the representatives to lie below $N^2$.
This ordering avoids circular witness choices.

Choose $A=X_W(W+1)>0$, which is odd because $W$ is even.
The recurrence gives $Y_W(W+1)\equiv W\pmod{2W}$, so
\[
 z=\frac{Y_W(W+1)/W-1}{2}\in\Z
\]
satisfies the odd frame. Choose all principal pairs canonically, with
$U=X_h(A),C=Y_h(A)$. If $k_*$ is the largest other requested index,
Pell doubling and growth give
\[
 Y_{2k_*+1}(A)>(4A-2)Y_{k_*}(A)^2.
\]
Together with~\eqref{eq:dominate314} and $A>W^W$, this implies
$C>W+\sum_{i\ne *}Y_{k_i}(A)^2$. The square-cap slack is therefore
nonnegative and has a ternary representation. The pooled quotient
exists because $Y_{k_i}(A)\equiv k_i\pmod{A-1}$.

For the auxiliary pair choose $l=hC$ and
\[
 H=X_l(A),\qquad J_c=Y_l(A)=C Y_C(U).
\]
The unit-power expansion gives an integer $T$ such that
\[
 Y_C(U)=CU^{C-1}+(U^2-1)T,
 \qquad
 \frac{J_c}{C^2}=U^{C-1}+(A^2-1)CT.
\]
The quotient is a positive odd integer, since $A,U$ are odd. Thus
$e=(J_c/C^2-1)/2$ supplies~\eqref{eq:oddaux314} for every semantic
index $h$. The even frame extends package 132's even-index odd-quotient
construction to this arbitrary anchor index.

Define $D$ by~\eqref{eq:detcenter314} and take
$S=X_h(D),Q=Y_h(D)$. These signed integers satisfy the comparison
norm, and polynomial base transfer through~\eqref{eq:centerresidues314}
gives both required congruences. Their unrestricted quotient witnesses
therefore exist. Fill the unchanged arithmetic gates, front guards and
power quotients. This constructs a zero of the complete new source.

The two existential directions prove the new arithmetic theorem. They
may replace old private Pell witnesses; no pointwise bijection fixing
those witnesses is asserted.

\subsection{Exact compilation and checked scope}

The offline rebuild starts at the pinned, literal, already factored parent
source. No new Horn factorization is performed. The source certificate
lists all twelve modified or removed rows and both appended center rows.
An independently written checker reconstructs the whole changed source,
including every frame atom and all fifteen pool summands.

Integral monic projection then restores all 216 eliminated coordinates
by exact integer polynomials, checks all 94 genuinely absent free
coordinates, and verifies the complete retained maps. The final 62
coordinate translations are integral bijections, substituted into every
row. Every coefficient of the sum of 204 squares is independently
collected. Both the final and parent metrics are computed from literal
equations. Three altered-artifact controls are rejected.

These checks import no producer, optimizer or translation helper. The
separate rebuild regenerates nine mathematical files byte for byte in a
temporary directory. They establish finite algebraic claims; the
unbounded Pell, representation, arithmetic and syntactic claims are
the ordinary proofs above and in the retained dependencies.
